\section{Methodology}
\paragraph{Notation.} We denote input sequences as $\mathbf{x} \in \mathbb{R}^{L \times F}$ with lookback $L$ and $F$ features, targets as $\mathbf{y} \in \mathbb{R}^{H \times F}$ with horizon $H$. Backbone predictions are $\hat{\mathbf{y}}$, refined predictions $\tilde{\mathbf{y}}$. Superscripts $(T), (S), (I)$ index trend, seasonal, and irregular components respectively.

\subsection{Overview}

\begin{algorithm}[h!]
\small
\caption{CTRL: LLM-Guided Residual Learning}
\label{alg:ctrl}
\begin{algorithmic}[1]
\REQUIRE Frozen backbone $f_\theta$, LLM agents, threshold $\tau{=}2.0$, check interval $K{=}50$
\ENSURE Refined prediction $\tilde{y} = \hat{y} + \delta$

\STATE \textcolor{blue}{\textit{// Stage 1: Control Signal Generation (no LLM parameter updates)}}
\STATE Compute global context: dataset-level error statistics
\STATE Cluster held-out training samples by temporal features via TypiClust
\STATE Retrieve $k{=}16$ few-shot examples with measured backbone errors
\STATE Compute current context: STL comparison of prediction vs ground truth
\FOR{agent $\in$ \{Trend, Seasonal\}}
    \STATE $z^{(\cdot)} \gets$ [scale, bias, gate, conf] via in-context reasoning
\ENDFOR
\STATE $z^{(I)} \gets$ encode Irregular agent text via frozen GPT-2
\STATE All samples share the same control signal $z$ \COMMENT{3 LLM calls total}

\STATE \textcolor{blue}{\textit{// Stage 2: Residual Decoder Training}}
\FOR{each $(x, y)$ in held-out training set}
    \STATE $\hat{y} \gets f_\theta(x)$
    \STATE $(\hat{T}, \hat{S}, \hat{I}) \gets \textsc{STL}(\hat{y})$
    \STATE $\delta^{(c)} \gets \text{MLP}_c([\hat{c}, z])$ for $c \in \{T, S, I\}$
    \STATE $\delta \gets w_T \delta^{(T)} + w_S \delta^{(S)} + w_I \delta^{(I)}$
    \STATE Minimize $\|\hat{y} + \delta - y\|^2$
\ENDFOR

\STATE \textcolor{blue}{\textit{// Stage 3: Test-Time Adaptation (every $K$ batches; Irregular bypasses)}}
\STATE Cache STL statistics from held-out training inputs
\FOR{each test batch}
    \STATE $\hat{y} \gets f_\theta(x_{\text{test}})$
    \STATE $s_c \gets$ z-score of $\textsc{STL}(x_{\text{test}})$ vs held-out training
    \IF{$s_c > \tau$ for $c \in \{T, S\}$}
        \STATE $z^{(c)} \gets$ LLM-adjusted control signal
    \ENDIF
    \STATE $\tilde{y} \gets \hat{y} + g_\phi(\hat{y}, z)$
\ENDFOR
\end{algorithmic}
\end{algorithm}

CTRL refines backbone forecasts through LLM-guided residual learning. LLMs excel at grounded reasoning, diagnosing discrepancies between predictions and ground truth, but struggle with precise numerical prediction. We exploit this asymmetry: LLM agents function as controllers that analyze prediction errors and generate compact policies guiding a learned residual decoder, while the backbone maintains numerical precision.

Given a multivariate time series with input $\mathbf{x} \in \mathbb{R}^{L \times F}$ and target $\mathbf{y} \in \mathbb{R}^{H \times F}$, a pretrained backbone $f_\theta$ produces initial predictions $\hat{\mathbf{y}} = f_\theta(\mathbf{x})$. CTRL refines these via a control signal-guided residual decoder:
\begin{equation}
    \tilde{\mathbf{y}} = \hat{\mathbf{y}} + g_\phi(\hat{\mathbf{y}}, \mathbf{z})
\end{equation}
where $g_\phi$ is the residual decoder with parameters $\phi$, and $\mathbf{z} \in \mathbb{R}^{d_z}$ is the control signal vector generated by LLM agents ($d_z = 8 + d_I$, where $d_I$ is the text projection dimension with default $d_I = 32$).

As shown in Figure \ref{fig:pipeline} and Algorithm~\ref{alg:ctrl}, CTRL operates in three stages: (1) control signal generation, where LLM agents produce correction policies via in-context reasoning without any LLM parameter updates; (2) residual decoder training on held-out training data; and (3) test-time adaptation that adjusts policies when distribution shift is detected.
Critically, test-time adaptation requires \emph{no ground-truth labels}, enabling adaptation to distribution shifts using only input statistics. LLM agents require no fine-tuning. Policies emerge from few-shot in-context reasoning over STL statistics. %
\subsection{Backbone Pretraining}

We train a time series backbone on training data, then freeze its weights. We evaluate DLinear~\citep{Zeng2023DLinear} (linear) and PatchTST~\citep{Nie2023PatchTST} (transformer), representing two dominant architectural paradigms in time series forecasting.
\subsection{Multi-Agent Control Signal Generation}

Control signal generation analyzes \emph{how} the backbone fails by comparing predictions against ground truth on held-out training data. 
Rather than using LLMs as numerical predictors, we position them as controllers that generate strategic correction policies while delegating computation to the residual decoder. Each agent receives multi-scale context to ground its reasoning.
This process requires no LLM parameter updates: agents perform in-context control signal generation by reasoning over held-out training statistics in a single forward pass. No gradients flow through the LLM at any stage; only the residual decoder is trained.

Effective control signal generation requires grounding at multiple scales. At the global level, agents receive dataset-wide statistics, such as mean error magnitudes, trend directions, seasonal strength distributions, establishing baseline expectations. Local context provides few-shot examples from similar held-out training samples, showing concrete failure patterns with measured errors that adapt correction magnitudes. 
Current context supplies STL decomposition comparing the sample's prediction against ground truth, including trend slope, seasonal amplitude, and dominance metrics (Appendix~\ref{app:stl_context}).


We select examples via TypiClust~\cite{hacohen2022active}, clustering held-out training samples by temporal features $\mathbf{f}_i = [\text{slope}, \sigma, \rho_1, S_{\text{strength}}]$ (trend slope, volatility, first-order autocorrelation, seasonal strength). We retrieve $k=16$ examples prioritizing same-cluster members while including diverse contrasts.

These held-out training examples include ground truth, enabling display of \emph{actual} backbone errors:
\begin{quote}
\small
\textit{``Similar pattern (upward trend, high volatility): Backbone MAE was 0.082, underestimating trend slope by 15\%.''}
\end{quote}
This grounds control signal decisions in empirical magnitudes rather than pattern inference alone.

We decompose both backbone prediction and ground truth into trend, seasonal, and irregular components using STL, providing agents with structured comparison of temporal characteristics.
STL decomposition partitions prediction errors into distinct frequency bands, each exhibiting characteristic failure modes that benefit from specialized analysis. The {Trend Agent} examines low-frequency dynamics including direction, level shifts, and systematic bias, detecting momentum underestimation when backbones predict flat trends against rising ground truth. The {Seasonal Agent} targets periodic patterns, diagnosing amplitude dampening and phase misalignment that backbones often introduce through oversmoothing. The {Irregular Agent} assesses high-frequency noise characteristics, prescribing adjustments when backbones either over-smooth or amplify fluctuations, enabling each agent to develop focused expertise. %

Each numerical agent (Trend, Seasonal) outputs a 4-dimensional control signal vector:
\begin{equation}
    \mathbf{z}^{(c)} = [\texttt{scale}, \texttt{bias}, \texttt{gate}, \texttt{confidence}] \in \mathbb{R}^4
\end{equation}
where $c \in \{T, S\}$ indexes trend and seasonal components, and \texttt{scale} $\in [0.5, 1.5]$ and \texttt{bias} $\in [-1, 1]$ provide linear modulation~\cite{perez2018film}, \texttt{gate} $\in [0, 1]$ controls correction aggressiveness inspired by gating mechanisms~\cite{hochreiter1997lstm}, and \texttt{confidence} $\in [0, 1]$ expresses agent certainty for uncertainty-aware weighting~\cite{xiong2024llms-conf,geng2024survey-conf}.

The Irregular Agent produces natural language analysis encoding uncertainty and anomaly patterns, projected to a $d_I$-dimensional embedding via a frozen text encoder:
\begin{equation}
    \mathbf{z}^{(I)} = \texttt{TextProj}(\texttt{GPT-2}(\mathbf{t}^{(I)})) \in \mathbb{R}^{d_I}
\end{equation}
where \texttt{TextProj} is a lightweight MLP projector that maps the 768D GPT-2 embedding to $d_I$ dimensions.

Prior work produces high-dimensional embeddings that downstream projections largely discard~\cite{Tan2024LLMUseful}. Our bounded output aligns with LLM strengths, structured values with clear semantics, while forcing distillation of actionable guidance.


\paragraph{Structural isolation from pretraining data.} LLM agents never receive raw time series values or dataset identifiers. Their inputs consist solely of backbone-specific runtime statistics: per-sample error magnitudes, slope deviations between ground truth and a specific backbone's predictions on a specific held-out split, seasonal amplitude mismatches, and dominance ratios. Concrete examples include phrases such as ``DLinear underestimates trend slope by 15\%'' or ``seasonal amplitude dampened from 0.14 to 0.09'', which are generated at runtime and depend on the backbone instance, the held-out split, and the STL decomposition parameters. These artifacts cannot exist in any pretraining corpus, making data leakage through LLM memorization architecturally precluded. This is supported empirically: removing STL context from prompts degrades performance by +2.11\%/+3.10\% (Table~\ref{tab:ablation}), confirming the LLM reasons over provided inputs rather than retrieving memorized patterns.



\begin{table*}[h!]
\centering
\footnotesize
\setlength{\tabcolsep}{1.1pt}
\newcolumntype{b}{>{\columncolor{cyan!5}}c}
\begin{tabular}{clbbbbcccccccccccc}
\toprule
\textbf{Data} & \textbf{H} & \multicolumn{2}{c}{\cellcolor{cyan!5}\textbf{Ours(DL)}} & \multicolumn{2}{c}{\cellcolor{cyan!5}\textbf{Ours(PT)}} & \multicolumn{2}{c}{\textbf{DLinear}} & \multicolumn{2}{c}{\textbf{PatchTST}} & \multicolumn{2}{c}{\textbf{TimeLLM$^\dagger$}} & \multicolumn{2}{c}{\textbf{CALF}} & \multicolumn{2}{c}{\textbf{TEMPO}} & \multicolumn{2}{c}{\textbf{GPT4TS}} \\
& & \cellcolor{cyan!5}\textbf{MAE} & \cellcolor{cyan!5}\textbf{MSE} & \cellcolor{cyan!5}\textbf{MAE} & \cellcolor{cyan!5}\textbf{MSE} & \textbf{MAE} & \textbf{MSE} & \textbf{MAE} & \textbf{MSE} & \textbf{MAE} & \textbf{MSE} & \textbf{MAE} & \textbf{MSE} & \textbf{MAE} & \textbf{MSE} & \textbf{MAE} & \textbf{MSE} \\
\midrule
\multirow{5}{*}{\rotatebox{90}{ETTh1}} 
& 48 & \textbf{.378} & \textbf{.340} & .381 & .342 & .382 & .346 & .386 & .346 & .398 & .369 & .400 & .372 & .419 & .395 & .425 & .412 \\
& 96 & \textbf{.398} & \textbf{.371} & .399 & .372 & .399 & .372 & .405 & .380 & .429 & .414 & .410 & .393 & .469 & .471 & .464 & .473 \\
& 192 & \textbf{.429} & \textbf{.410} & .439 & .429 & .430 & .412 & .439 & .431 & .432 & .425 & .436 & .430 & .494 & .514 & .527 & .598 \\
& 336 & .457 & .442 & \textbf{.430} & \textbf{.415} & .459 & .444 & .431 & .416 & .452 & .450 & .476 & .490 & .477 & .485 & .599 & .738 \\
& 720 & .510 & .498 & \textbf{.458} & \textbf{.443} & .511 & .497 & .459 & .444 & .477 & .462 & .522 & .540 & .502 & .502 & .600 & .783 \\
\midrule
\multirow{5}{*}{\rotatebox{90}{ETTh2}} 
& 48 & \textbf{.296} & \textbf{.217} & .309 & .242 & .318 & .235 & .309 & .242 & .318 & .243 & .316 & .243 & .334 & .263 & .343 & .280 \\
& 96 & \textbf{.335} & \textbf{.273} & .356 & .305 & .370 & .307 & .410 & .343 & .359 & .307 & .367 & .320 & .379 & .332 & .344 & .284 \\
& 192 & \textbf{.388} & \textbf{.349} & .408 & .373 & .427 & .395 & .427 & .376 & .392 & .357 & .421 & .409 & .406 & .378 & .421 & .412 \\
& 336 & .427 & .383 & \textbf{.406} & \textbf{.355} & .455 & .438 & .414 & .364 & .412 & .375 & .421 & .380 & .422 & .399 & .442 & .429 \\
& 720 & .533 & .501 & .465 & \textbf{.408} & .575 & .650 & .482 & .458 & .463 & .442 & \textbf{.442} & .416 & .462 & .436 & .477 & .475 \\
\midrule
\multirow{5}{*}{\rotatebox{90}{ETTm1}} 
& 48 & .329 & .272 & \textbf{.319} & .264 & .332 & .276 & \textbf{.319} & .268 & .324 & \textbf{.261} & .325 & .266 & .358 & .302 & .394 & .383 \\
& 96 & \textbf{.344} & \textbf{.297} & .382 & .362 & .347 & .301 & .393 & .404 & .350 & \textbf{.297} & .347 & .298 & .382 & .337 & .395 & .384 \\
& 192 & \textbf{.366} & .330 & .427 & .439 & .368 & .338 & .436 & .473 & .372 & .333 & .374 & \textbf{.320} & .408 & .390 & .424 & .423 \\
& 336 & \textbf{.388} & \textbf{.370} & .414 & .409 & .390 & .371 & .419 & .417 & .398 & .379 & .394 & .371 & .434 & .421 & .444 & .458 \\
& 720 & \textbf{.421} & .421 & .446 & .446 & .429 & .430 & .454 & .460 & .425 & \textbf{.415} & .427 & .428 & .462 & .436 & .479 & .525 \\
\midrule
\multirow{5}{*}{\rotatebox{90}{ETTm2}} 
& 48 & .226 & \textbf{.126} & \textbf{.222} & .129 & .230 & .129 & .223 & .130 & .236 & .136 & .226 & .128 & .240 & .146 & .245 & .151 \\
& 96 & \textbf{.251} & \textbf{.161} & .278 & .196 & .271 & .173 & .278 & .197 & .267 & .177 & .258 & .170 & .272 & .190 & .266 & .187 \\
& 192 & \textbf{.292} & \textbf{.214} & .315 & .247 & .311 & .230 & .319 & .256 & .309 & .238 & .296 & .226 & .320 & .271 & .315 & .260 \\
& 336 & \textbf{.328} & \textbf{.265} & .363 & .331 & .350 & .287 & .374 & .367 & .341 & .290 & .330 & .278 & .357 & .324 & .350 & .317 \\
& 720 & .386 & \textbf{.354} & .418 & .422 & .412 & .386 & .436 & .478 & .388 & .368 & \textbf{.382} & .362 & .413 & .409 & .406 & .408 \\
\midrule
\multirow{5}{*}{\rotatebox{90}{Weather}} 
& 48 & .190 & .135 & \textbf{.159} & \textbf{.118} & .196 & .137 & \textbf{.159} & .119 & .178 & .127 & .160 & .123 & .199 & .146 & .176 & .129 \\
& 96 & .225 & .169 & \textbf{.201} & \textbf{.156} & .232 & .172 & .209 & .158 & .212 & .161 & \textbf{.201} & .158 & .210 & \textbf{.156} & .217 & .167 \\
& 192 & .256 & \textbf{.200} & .265 & .232 & .276 & .216 & .275 & .243 & .250 & .202 & \textbf{.241} & \textbf{.200} & .259 & .206 & .259 & .216 \\
& 336 & .297 & .258 & .311 & .278 & .310 & .261 & .315 & .304 & .289 & .252 & \textbf{.282} & \textbf{.250} & .351 & .394 & .297 & .266 \\
& 720 & .343 & \textbf{.313} & .360 & .330 & .358 & .321 & .365 & .373 & .338 & .323 & \textbf{.330} & .320 & .347 & .339 & .343 & .330 \\
\midrule
\multirow{5}{*}{\rotatebox{90}{ECL}} 
& 48 & .226 & .127 & .205 & \textbf{.108} & .233 & .130 & .205 & .111 & .218 & .116 & \textbf{.200} & .109 & .212 & .114 & .234 & .131 \\
& 96 & .241 & .142 & \textbf{.224} & \textbf{.131} & .247 & .145 & \textbf{.224} & \textbf{.131} & .231 & .132 & .227 & .136 & .233 & .134 & .247 & .147 \\
& 192 & .254 & .155 & \textbf{.242} & \textbf{.149} & .259 & .158 & .243 & .150 & .250 & .150 & .245 & .156 & .246 & .152 & .258 & .162 \\
& 336 & .270 & .170 & \textbf{.258} & \textbf{.165} & .272 & .171 & .259 & \textbf{.165} & .266 & .167 & .263 & .171 & .264 & .171 & .275 & .179 \\
& 720 & .303 & .204 & .292 & .201 & .303 & .204 & .292 & .202 & .296 & .202 & \textbf{.276} & \textbf{.190} & .294 & .204 & .316 & .232 \\
\midrule
\multirow{5}{*}{\rotatebox{90}{Exch.}} 
& 24 & .104 & \textbf{.024} & \textbf{.101} & \textbf{.024} & .105 & \textbf{.024} & .105 & .025 & .135 & .036 & .106 & .025 & .109 & .026 & .118 & .029 \\
& 48 & .147 & \textbf{.044} & \textbf{.146} & \textbf{.044} & .150 & .045 & .149 & .048 & .164 & .053 & \textbf{.146} & \textbf{.044} & .154 & .049 & .157 & .050 \\
& 96 & \textbf{.201} & \textbf{.080} & \textbf{.201} & .083 & .208 & .082 & .212 & .094 & .234 & .107 & .206 & .087 & .207 & .088 & .249 & .121 \\
& 192 & .305 & \textbf{.168} & .308 & .172 & .325 & .186 & .311 & .196 & .349 & .229 & \textbf{.302} & .181 & .310 & .192 & .324 & .203 \\
& 336 & .418 & \textbf{.311} & \textbf{.416} & .334 & .422 & .317 & .418 & .334 & .476 & .409 & .429 & .353 & .443 & .368 & .454 & .383 \\
\bottomrule
\end{tabular}
\caption{Forecasting results on multivariate benchmarks. \textbf{Bold}: best. (DL) and (PT) indicate DLinear and PatchTST backbones, respectively. $^\dagger$Time-LLM~\cite{Jin2024TimeLLM} uses GPT-2 backend with a single seed; all other results averaged over 3 seeds.}
\label{tab:main_results}
\end{table*}


\subsection{Residual Decoder Training}

With frozen backbone and frozen LLM, we train a lightweight residual decoder on held-out training data. The decoder (${\sim}$400K parameters) is the only component updated via gradient descent, making it orders of magnitude smaller than the frozen modules it conditions on.




The decoder maintains component separation. We  apply STL decomposition to backbone predictions into trend, seasonal, and irregular components:
\begin{equation}
    \hat{\mathbf{y}} = \hat{\mathbf{T}} + \hat{\mathbf{S}} + \hat{\mathbf{I}}
\end{equation}
Three parallel MLPs process each component, receiving the concatenated control signal:
\begin{align}
    \boldsymbol{\delta}^{(T)} &= \texttt{MLP}_T([\hat{\mathbf{T}}, \mathbf{z}]) \\
    \boldsymbol{\delta}^{(S)} &= \texttt{MLP}_S([\hat{\mathbf{S}}, \mathbf{z}]) \\
    \boldsymbol{\delta}^{(I)} &= \texttt{MLP}_I([\hat{\mathbf{I}}, \mathbf{z}])
\end{align}
This enforces architectural alignment between specialized agents and their target components.


While control signal parameters are concatenated to MLP inputs rather than architecturally enforced, end-to-end training encourages the decoder to learn appropriate mappings. 

The residual decoder minimizes MSE between corrected predictions and ground truth:
\begin{equation}
    \mathcal{L} = \frac{1}{N}\sum_{i=1}^{N} \lVert \hat{\mathbf{y}}_i + \boldsymbol{\delta}_i - \mathbf{y}_i \rVert^2
\end{equation}
where the combined correction $\boldsymbol{\delta} = w_T \boldsymbol{\delta}^{(T)} + w_S \boldsymbol{\delta}^{(S)} + w_I \boldsymbol{\delta}^{(I)}$ aggregates component-wise outputs with learned weights $w_T, w_S, w_I$. 


\subsection{Test-Time Adaptation}

Policies trained on held-out training must generalize to potentially shifted test distributions without ground truth labels. We address this via label-free adaptation that detects distribution shift from input statistics alone and adjusts control signal parameters accordingly.

During control signal generation, we cache STL statistics from held-out training inputs: trend slope distribution, seasonal amplitude range, and irregular variance. At test time, we compute per-component divergence between test and held-out training statistics:
\begin{align}
    s_T &= \frac{|\text{slope}_\text{test} - \bar{\text{slope}}_\text{val}|}{\sigma_{\text{slope,val}} + \epsilon} \\
    s_S &= \frac{|d_{S,\text{test}} - \bar{d}_{S,\text{val}}|}{\sigma_{d_S,\text{val}} + \epsilon}
\end{align}
where $\bar{\cdot}$ and $\sigma$ denote held-out training mean and standard deviation respectively. This z-score formulation yields an adaptive threshold: we trigger adaptation when $s_c > \tau$ with $\tau = 2.0$, corresponding to test statistics outside two standard deviations of the held-out training distribution.

When adaptation triggers for component $c \in \{T, S\}$, the LLM agent receives shift statistics and outputs adjustment factors that modulate the original control signal through the same four-parameter interface used in control signal generation:
\begin{align}
    \texttt{scale}' &= \texttt{scale} \times m_\text{scale} \\
    \texttt{bias}' &= \texttt{bias} + \Delta_\text{bias} \\
    \texttt{gate}' &= \texttt{gate} \times m_\text{gate} \\
    \texttt{conf}' &= \texttt{conf} \times (1 - p_\text{conf})
\end{align}
where multipliers $m_* \in [0.5, 1.5]$ scale proportionally to shift severity, and confidence penalty $p_\text{conf} \in [0, 0.5]$ reflects increased uncertainty under distribution shift. Unshifted components retain original policies. The irregular agent uses text embedding. At inference, the adapted control signal $\mathbf{z}'$ passes through the same residual decoder interface, ensuring architectural consistency between training and test-time adaptation.

CTRL requires 3 LLM calls for initial control signal generation (one per agent), 
with test-time adaptation adding at most 3 calls per calibration check. This yields 3--24 total calls across datasets.







